%% Hoja de vida maestra (ESPAÑOL) - Juan Diego Chaves Guerrero
%% Template: minimal-word (templates/cv/minimal-word). One page, single column, ATS-safe.
%% Spanish master: /apply and /postular start from this file for Spanish postings,
%% and from cv/main_example.tex for English postings. Keep both masters in sync.
%% CONFIDENTIAL: never name the internal AI project built at Periferia IT Group.
%% Compile with: lualatex main_example_es.tex (preamble.tex must be in the same folder)
% Shared preamble for the minimal-word CV template (single column, black text,
% bold/italic hierarchy, Calibri-metric font). ATS-safe: no columns, no icons,
% contact details as literal text.
\documentclass[10pt,letterpaper]{article}
\usepackage[margin=0.55in,top=0.45in,bottom=0.45in]{geometry}
\usepackage{fontspec}
% Calibri on Windows (as in the candidate's Word CV); Carlito, its metric clone, elsewhere.
\IfFontExistsTF{Calibri}{\setmainfont{Calibri}}{\setmainfont{Carlito}}
\usepackage{xcolor}
\definecolor{meta}{gray}{0.40}        % grey for secondary text only; no colour accents
\usepackage{enumitem}
\usepackage{titlesec}
\usepackage{needspace}
\usepackage[hidelinks]{hyperref}
\pagestyle{empty}
\setlength{\parindent}{0pt}
\setlength{\parskip}{0pt}
\linespread{1.0}

% Section heading: bold small caps-style uppercase with a rule underneath
\titleformat{\section}{\normalsize\bfseries}{}{0pt}{\MakeUppercase}[\vspace{-8pt}\rule{\linewidth}{0.6pt}]
\titlespacing*{\section}{0pt}{5pt}{1pt}

\setlist[itemize]{leftmargin=1.6em,itemsep=0pt,topsep=1pt,parsep=0pt,label=\textbullet}

% \entry{Bold title}{grey meta}{right-aligned date}{italic subtitle}
\newcommand{\entry}[4]{%
  \needspace{3\baselineskip}%
  \textbf{#1}\ifx&#2&\else\ \textcolor{meta}{$|$\ #2}\fi\hfill\textcolor{meta}{#3}\par
  \ifx&#4&\else{\itshape #4}\par\fi}

% \skill{Label}{items}
\newcommand{\skill}[2]{\textbf{#1:}\ #2\par\vspace{1pt}}
% ATS: LuaTeX writes no space characters, so pypdf glues words together unless the
% fonts carry a real space glyph. \atsspaces embeds one per font, zero width; keep it
% at the start of the name line.
\newcommand{\atsspaces}{\makebox[0pt]{{\normalfont\char32}{\normalfont\itshape\char32}{\normalfont\bfseries\char32}}}

\ifpdftex\usepackage[T1]{fontenc}\fi % no-op under lualatex (this template needs fontspec)
\hypersetup{pdftitle={Juan Diego Chaves Guerrero - Hoja de vida}, pdflang={es-CO}}
\begin{document}

\begin{center}
{\large\bfseries\atsspaces Juan Diego Chaves Guerrero}\\[1pt]
{\small +57 316 258 8234\ \ $|$\ \ juand.chavesg@gmail.com\ \ $|$\ \ \href{https://linkedin.com/in/Juan-Diego-Chaves-guerrero}{linkedin.com/in/Juan-Diego-Chaves-guerrero}\ \ $|$\ \ \href{https://github.com/Juand2021}{github.com/Juand2021}}\\
{\small Bogotá, Colombia}\\[1pt]
{\small\itshape\textcolor{meta}{Ingeniero de IA\ \ $\cdot$\ \ Ingeniería de Sistemas}}
\end{center}
\vspace{-8pt}

\section{Perfil}
Ingeniero de IA enfocado en desarrollo agéntico, RAG y automatización con IA, con soluciones de IA para equipos internos en la Torre de IA de Periferia IT Group. Usuario avanzado de Claude Code; construye flujos de agentes con LangGraph y conecta agentes a sistemas reales mediante MCP. Ingeniero de Sistemas (Universidad Jorge Tadeo Lozano) con tres certificaciones Microsoft en IA/ML.

\section{Competencias Clave}
\skill{IA Generativa y LLMs}{Generación Aumentada por Recuperación (RAG), agentes basados en grafos con LangGraph, bases de datos vectoriales, indexación para recuperación rápida, orquestación de subagentes}
\skill{Integración de Agentes}{Integraciones con Model Context Protocol (MCP) y servidores MCP propios, automatización de procesos de negocio con IA, ingeniería de prompts}
\skill{Agentes de Código con IA}{Claude Code (avanzado), Cursor, Google Antigravity}
\skill{Programación}{Python (avanzado), TypeScript, Node.js, Electron, Java, Kotlin, C\#, JavaScript}
\skill{Datos y ML}{Supabase, Google Firestore, SQL Server; Pandas, Scikit-learn, aprendizaje supervisado, minería de datos}
\skill{Nube y Metodologías}{Microsoft Azure Machine Learning (fundamentos certificados), Git / GitHub, Scrum / Kanban}

\section{Experiencia Laboral}
\entry{Periferia IT Group}{Bogotá, Colombia (híbrido)}{Jun 2026 -- Dic 2026}{Practicante de Ingeniería de IA -- Torre de IA}
\begin{itemize}
  \item Construí soluciones y automatizaciones de IA para contabilidad, software y talento humano, y acompañé su implementación.
  \item Contribuí al desarrollo de un proyecto interno de IA, construyendo automatizaciones sobre él.
  \item Desarrollo de software agéntico con Claude Code, Cursor y Google Antigravity, orquestando subagentes y prompts estructurados.
  \item Conecté agentes de IA a Supabase, bases de datos y herramientas frontend/backend vía MCP; construí servidores MCP propios.
  \item Stack: Python, TypeScript, Node.js, aplicaciones de escritorio con Electron, Supabase, Google Firestore.
\end{itemize}

\entry{Soluciones Empresariales 360° (SOE 360)}{Bogotá, Colombia}{2022 $\cdot$ 6 meses}{Analista de Mesa de Servicio TI -- Practicante}
\begin{itemize}
  \item Adopté ChatGPT en sus primeros meses, usándolo para escribir código y ajustar soluciones de TI para la mesa de servicio.
  \item Resolví incidentes bajo prácticas ITIL y documenté soluciones recurrentes, reduciendo los tiempos de resolución.
\end{itemize}

\section{Proyectos}
\entry{HeartAttackPrevention -- Predictor de Riesgo Cardíaco}{Proyecto Personal}{2025}{Python $\cdot$ Scikit-learn $\cdot$ Árbol de Decisión $\cdot$ GridSearchCV}
\begin{itemize}
  \item App web de riesgo de infarto (85,5\,\% de precisión); ganadora de la Feria de Ingeniería de la Universidad Jorge Tadeo Lozano.
\end{itemize}
\entry{Planorama -- Recomendador Inteligente de Eventos}{Proyecto Personal}{2024}{Python $\cdot$ Streamlit $\cdot$ TF-IDF $\cdot$ Pandas $\cdot$ Folium}
\begin{itemize}
  \item Recomendador conversacional de eventos en Bogotá; resultados relevantes en menos de 2 segundos y mapa geolocalizado.
\end{itemize}
\entry{TadIA -- Asistente Académico Android}{Proyecto Personal}{2026}{Android $\cdot$ Kotlin $\cdot$ Jetpack Compose}
\begin{itemize}
  \item Aplicación Android para estudiantes universitarios con autenticación y módulos académicos, construida con Material Design 3.
\end{itemize}

\section{Educación}
\entry{Universidad Jorge Tadeo Lozano}{Bogotá, Colombia}{2023 -- 2026 $\cdot$ Culminado}{Ingeniería de Sistemas -- materias y práctica culminadas; grado en febrero de 2027}
\entry{SENA -- Servicio Nacional de Aprendizaje}{Bogotá, Colombia}{Graduado 2023}{Tecnólogo en Análisis y Desarrollo de Sistemas de Información\ \ $\cdot$\ \ Técnico en Programación de Software}

\section{Certificaciones}
\entry{Building Intelligent Troubleshooting Agents}{Microsoft}{Nov 2026}{}
\entry{AI and Machine Learning Algorithms and Techniques}{Microsoft}{Nov 2026}{}
\entry{Foundations of AI and Machine Learning}{Microsoft}{Nov 2026}{}

\section{Idiomas y Reconocimientos}
Español (nativo)\ \ $\cdot$\ \ Inglés (B2)\ \ $\cdot$\ \ Ganador, Feria de Ingeniería de la Universidad Jorge Tadeo Lozano (2025)

\end{document}
